% Options for packages loaded elsewhere
\PassOptionsToPackage{unicode}{hyperref}
\PassOptionsToPackage{hyphens}{url}
\PassOptionsToPackage{dvipsnames,svgnames,x11names}{xcolor}
%
\documentclass[
  11pt,
]{article}
\usepackage{amsmath,amssymb}
\usepackage{iftex}
\ifPDFTeX
  \usepackage[T1]{fontenc}
  \usepackage[utf8]{inputenc}
  \usepackage{textcomp} % provide euro and other symbols
\else % if luatex or xetex
  \usepackage{unicode-math} % this also loads fontspec
  \defaultfontfeatures{Scale=MatchLowercase}
  \defaultfontfeatures[\rmfamily]{Ligatures=TeX,Scale=1}
\fi
\usepackage{lmodern}
\ifPDFTeX\else
  % xetex/luatex font selection
\fi
% Use upquote if available, for straight quotes in verbatim environments
\IfFileExists{upquote.sty}{\usepackage{upquote}}{}
\IfFileExists{microtype.sty}{% use microtype if available
  \usepackage[]{microtype}
  \UseMicrotypeSet[protrusion]{basicmath} % disable protrusion for tt fonts
}{}
\makeatletter
\@ifundefined{KOMAClassName}{% if non-KOMA class
  \IfFileExists{parskip.sty}{%
    \usepackage{parskip}
  }{% else
    \setlength{\parindent}{0pt}
    \setlength{\parskip}{6pt plus 2pt minus 1pt}}
}{% if KOMA class
  \KOMAoptions{parskip=half}}
\makeatother
\usepackage{xcolor}
\usepackage[margin=1in]{geometry}
\usepackage{longtable,booktabs,array}
\usepackage{calc} % for calculating minipage widths
% Correct order of tables after \paragraph or \subparagraph
\usepackage{etoolbox}
\makeatletter
\patchcmd\longtable{\par}{\if@noskipsec\mbox{}\fi\par}{}{}
\makeatother
% Allow footnotes in longtable head/foot
\IfFileExists{footnotehyper.sty}{\usepackage{footnotehyper}}{\usepackage{footnote}}
\makesavenoteenv{longtable}
\setlength{\emergencystretch}{3em} % prevent overfull lines
\providecommand{\tightlist}{%
  \setlength{\itemsep}{0pt}\setlength{\parskip}{0pt}}
\setcounter{secnumdepth}{-\maxdimen} % remove section numbering
\ifLuaTeX
  \usepackage{selnolig}  % disable illegal ligatures
\fi
\usepackage{bookmark}
\IfFileExists{xurl.sty}{\usepackage{xurl}}{} % add URL line breaks if available
\urlstyle{same}
\hypersetup{
  pdftitle={How large can the restricted order of an additive basis be?},
  pdfauthor={Leo Gao},
  colorlinks=true,
  linkcolor={Maroon},
  filecolor={Maroon},
  citecolor={Blue},
  urlcolor={teal},
  pdfcreator={LaTeX via pandoc}}

\title{How large can the restricted order of an additive basis be?}
\author{Leo Gao}
\date{27 November 2026 · AI-assisted research manuscript}

\begin{document}
\maketitle

\textbf{Abstract.} Erdős and Graham asked whether there is a function of
ordinary order bounding the finite restricted order of an asymptotic
additive basis. We give an affirmative answer. More precisely, let
\(F(h)\) be the maximum finite restricted order among bases of ordinary
order exactly \(h\). We prove \(\log\log F(h)=\Theta(h)\) as
\(h\to\infty\). The upper bound combines a finite-group lemma on
disjoint subset sums with quantitative inputs of Hegyvári, Hennecart and
Plagne and of Chen, Yang and Zhao. For the lower bound, an eventually
periodic mixed-radix basis has ordinary order \(h\) but requires at
least \(2^{2^{h-5}+h-4}\) distinct summands for \(h\ge5\). The problem
of characterizing when ordinary and restricted orders are equal is not
addressed here.

\subsection{1. The extremal question}\label{the-extremal-question}

For \(A\subseteq\mathbb N_0\), write \(S_j(A)\) for sums of at most
\(j\) members of \(A\), with repetition allowed, and \(R_j(A)\) for sums
of at most \(j\) distinct members. Both sets include \(0\). The ordinary
order \(\operatorname{ord}(A)\), respectively restricted order
\(\operatorname{ord}_R(A)\), is the least \(j\ge1\) for which
\(S_j(A)\), respectively \(R_j(A)\), is cofinite; a nonexistent order is
infinite.

Erdős and Graham {[}1, p.~206; 2, p.~52{]} asked whether the finite
restricted order can be bounded in terms of the ordinary order.
Hegyvári, Hennecart and Plagne {[}4, Theorem 4{]} constructed bases of
order \(h\ge3\) with restricted order at least \(2^{h-2}+h-1\). Chen,
Yang and Zhao {[}5{]} have recently proved two related results: the
ordinary finite-deletion hypothesis implies finite restricted order, and
a basis has finite restricted order if and only if its finite subset
sums cover every residue class modulo every positive integer. Those
results give the existence criterion. The question here asks instead for
a bound uniform across different bases of the same ordinary order.

\textbf{Theorem 1.} There is an effectively specified function \(f\)
such that \(\operatorname{ord}_R(A)\le f(h)\) whenever
\(\operatorname{ord}(A)\le h\) and \(\operatorname{ord}_R(A)<\infty\).
If \(F(h)\) denotes the largest finite restricted order among bases of
ordinary order exactly \(h\), then for \(h\ge5\) \[
2^{\,2^{h-5}+h-4}\ \le\ F(h)\ \le\ \exp\!\bigl(O(16^h h\log(2h))\bigr).
\] In particular \(\log\log F(h)=\Theta(h)\). The lower-bound examples
remain bases after every finite deletion.

The point of the upper bound is to select a single budget for residue
corrections before the modulus that requires those corrections is known.

\subsection{2. A finite-group lemma}\label{a-finite-group-lemma}

Let \(I\) be a set of labels, each assigned a value in a finite abelian
group \(G\) of order \(q\). Different labels may have the same value,
and \(I\) need not be finite. Write \(R_j\) for sums of at most \(j\)
distinct labels, with \(R_0=\{0\}\), and \(\Sigma=\bigcup_{j\ge0}R_j\).

\textbf{Lemma 2.} Given \(h\ge1\) and \(0<\delta<1\), there is a number
\(b=b(h,\delta)\) independent of \(G\) and \(I\) such that \[
|R_h|\ge\delta q,\qquad \Sigma=G\qquad\Longrightarrow\qquad R_b=G.
\] One choice is given by \(\alpha_h=\delta\) and the downward recursion
\[
L_r=\left\lceil4\alpha_r^{-2}\right\rceil,\qquad
Q_r=(2r+1)L_r+1,\qquad
\alpha_{r-1}=\frac{\alpha_r^2}{16rQ_r}\quad(r=h,\ldots,1),
\] with
\(b=\max\{\lceil1/\alpha_0\rceil,\max_{1\le r\le h}(2r+1)L_r\}\).

\emph{Proof.} In any group of order \(g\), a shortest subset
representation of a reachable value uses at most \(g-1\) labels.
Otherwise two of its ordered partial sums, counting the initial zero,
coincide, allowing deletion of a zero-sum subsequence. This also applies
in every quotient. It proves the lemma for \(q\le1/\alpha_0\).

For larger \(q\), choose the least \(r\ge1\) such that
\(|R_r|\ge\alpha_r q\). Let \(X=R_r\), \(M=|X|\), and \(e=|R_{r-1}|\).
Then \(e\le\alpha_{r-1}q\), by minimality when \(r>1\) and by \(e=1\)
when \(r=1\). For \(x\in X\) fix one set \(U_x\) of at most \(r\) labels
representing \(x\), using \(U_0=\varnothing\). Any one label of value
\(v\) can occur in at most \(e\) of these sets: the corresponding
\(x-v\) are distinct members of \(R_{r-1}\).

Let \(D\) consist of \(0\) and those residues with \(Q_r\) pairwise
label-disjoint nonempty representations of length at most \(2r\). For
\(z\in G\) let \(p_z\) count ordered pairs \((x,y)\in X^2\) with
\(x+y=z\) and \(U_x\cap U_y=\varnothing\). At most \(rMe\) pairs fail
this disjointness condition, and \(p_z\le M\). For \(z\notin D\), take a
maximal family of counted pairs with mutually disjoint unions
\(U_x\cup U_y\). Each union is nonempty since \(z\ne0\). The union \(W\)
of the whole family has at most \(2r(Q_r-1)\) labels and hits every
counted pair. Each label in \(W\) can occur on the first side of at most
\(e\) pairs of sum \(z\) (fixing \(x\) fixes \(y\)), and likewise on the
second side. Thus \(p_z\le4r(Q_r-1)e\) for \(z\notin D\). Consequently
\[
M^2-rMe\le\sum_zp_z\le |D|q+4r(Q_r-1)eq,\qquad
\frac{|D|}{q}\ge \alpha_r^2-(r+4r(Q_r-1))\alpha_{r-1}
\ge\frac{\alpha_r^2}{2}.
\]

Let \(J=\langle D\rangle\). Its index in \(G\) is at most
\(2/\alpha_r^2\). We have \(L_rD=J\), where this sumset allows
repetitions. Indeed, if \(L_rD\ne J\), let \(U\) be its stabilizer
inside \(J\). Then \(D\) meets \(k\ge2\) cosets of \(U\) since
\(0\in D\) and \(D\) generates \(J\). By the finite abelian addition
theorem of Kneser {[}6{]}, \[
|L_rD|\ \ge\ (L_r(k-1)+1)|U|\ >\ (L_r/2)|D|\ \ge\ q,
\] which is impossible. The strict inequality uses \(k-1\ge k/2\) and
\(|D|\le k|U|\).

To represent \(z\in G\), choose a set \(E\) of labels representing its
coset in \(G/J\), with \(|E|\le [G:J]-1\le L_r-1\), using \(\Sigma=G\)
and the partial-sum observation. The difference between \(z\) and the
sum of \(E\) lies in \(J=L_rD\). Express it as \(L_r\) terms of \(D\).
For each nonzero term in turn, choose one of its \(Q_r\) disjoint
representations avoiding \(E\) and all labels used so far. There are at
most \((L_r-1)+2r(L_r-1)<Q_r\) forbidden labels; the zero term needs
none. The resulting representation uses at most \((2r+1)L_r\le b\)
distinct labels. \(\square\)

\subsection{3. Transfer from the quotients to the
integers}\label{transfer-from-the-quotients-to-the-integers}

Write \(\underline d\) for lower asymptotic density. Hegyvári, Hennecart
and Plagne {[}3; 7, Theorem 4.8{]} give \[
\underline d(R_h(A))\ge c_h,\qquad
c_h=h^{-h}\exp(-\pi\sqrt{2h/3})
\] if \(\operatorname{ord}(A)\le h\). To see the convention explicitly,
apply their exact-\(h\) density inequality to \(B=(A\cup\{0\})+1\). The
ordinary sumset \(hB\) is cofinite. An exact distinct \(h\)-sum from
\(B\), translated by \(-h\), is a distinct sum from \(A\cup\{0\}\); zero
may be omitted.

Set \(\delta=c_h/2\). For every \(q\ge1\), the residues of \(R_h(A)\)
cover at least \(\delta q\) classes, since a set contained in \(t\)
classes has upper density at most \(t/q\). If
\(\operatorname{ord}_R(A)<\infty\), its arbitrary finite subset sums
occupy every class. Apply Lemma 2 with \(I=A\), assigning to each
element its value modulo \(q\). Then \[
R_b(A)\bmod q=\mathbb Z/q\mathbb Z\qquad\text{for every }q\ge1,
\qquad b=b(h,\delta).
\] The same conclusion holds if the modular completeness condition alone
is assumed.

We use the quantitative content of Chen, Yang and Zhao {[}5, Lemma 4,
proof, eq. (2){]}. Given \(\underline d(R_h(A))>\delta\) and the above
budget \(b\), set \(\beta_h=\delta\), and for \(r=h,\dots,1\) put \[
T_r=\lceil 8\beta_r^{-2}\rceil,\quad
V_r=b+\sum_{j=r+1}^h 2jT_j+2r(T_r-1)+1,\quad
\beta_{r-1}=\frac{\beta_r^2}{64rV_r},\quad K=\sum_{j=1}^h2jT_j.
\] Their lemma gives integers \(g,N_0\ge1\) and a residue \(a\bmod g\)
such that all \(n\ge N_0\) with \(n\equiv a\bmod g\) lie in
\(R_K(A\setminus E)\) for \textbf{every} \(E\subseteq A\) with
\(|E|\le b\). The period \(g\) and the onset \(N_0\) can depend on
\(A\); the displayed \(K\) cannot.

For each residue \(v\bmod g\), use Lemma 2 to choose \(E_v\subseteq A\)
of size at most \(b\) with sum \(e_v\equiv v-a\bmod g\). Given
\(n\ge N_0+\max_v e_v\), set \(v=n\bmod g\). Then \(n-e_v\) belongs to
the progression and is represented by at most \(K\) distinct members
outside \(E_v\). Adjoining \(E_v\) proves the universal bound
\(f(h)=K+b\).

For its growth, put \(x_r=1/\alpha_r\). From \(L_r\le5x_r^2\) and
\(Q_r\le16h x_r^2\) we have \(x_{r-1}\le256h^2x_r^4\). It follows by
iteration that \[
\log b=O\bigl(4^h(\log(1/\delta)+\log(2h))\bigr).
\] Similarly \(y_r=1/\beta_r\) increases as \(r\) decreases and
\(T_j\le9y_r^2\) for \(j\ge r\). Thus \(V_r\le b+19h^2y_r^2+1\), whence
\(y_{r-1}\le Ch^3(b+1)y_r^4\) for some absolute \(C\). Iteration gives
\[
\log(K+b)=O\bigl(16^h(\log(1/\delta)+\log(2h))\bigr)
=O\bigl(16^h h\log(2h)\bigr).
\]

\subsection{4. A periodic mixed-radix
construction}\label{a-periodic-mixed-radix-construction}

Fix \(h\ge5\) and put \(d=h-4\). Define \[
P_0=1,\qquad P_i=2^{\,2^{i-1}+i}\ (1\le i\le d),\qquad
P=P_d,\qquad u_i=P_i/P_{i-1},\qquad w_i=P/P_i.
\] Index \(1\) denotes the highest digit. The radices are integers:
\(u_1=4\) and \(u_i=2^{2^{i-2}+1}\) for \(i\ge2\). Take the
denominations \[
C=\{P\}\cup\{kw_i:1\le i\le d,\ 1\le k<u_i\}.
\] Set \(m=8P\) and form the integer basis \[
A=\{2c:c\in C\}\ \cup\ \{mn,mn+1:n\ge1\}.
\] The finite exceptions are below \(m\). A number \(0\le t<4P\) is a
sum of at most three copies of \(P\) and one denomination for each
nonzero mixed digit. For a sufficiently large target congruent to
\(2t+\epsilon\bmod m\), with \(\epsilon\in\{0,1\}\), take those doubled
denominations and one tail term in residue \(\epsilon\). This proves
\(\operatorname{ord}(A)\le 1+3+d=h\).

To prove equality, consider \(n\equiv m-1\bmod m\). A representation
contains an odd number \(2v+1\) of tail terms of residue \(1\). At the
level of residues, discard all residue-zero terms, keep one \(1\), and
replace the other \(2v\) ones by \(v\) copies of the smallest doubled
denomination, \(2\). This cannot increase length. After subtracting the
remaining \(1\) and dividing the even sum by \(2\), we obtain a
nonnegative \(t\equiv4P-1\bmod 4P\) in denominations \(C\). Any
representation of \(t\) uses at least \(\lfloor t/P\rfloor+s(t\bmod P)\)
coins, where \(s\) counts nonzero digits. Indeed, normalize from low to
high. At a place of radix \(u\), the digits of \(c>0\) coins sum to
\(ku+v'\) with \(k\le c-1\) and \(0\le v'<u\). Replace them by \(k\)
units at the next place and at most one digit at the current place
without increasing length. The normalized length has the asserted value.
Now \(\lfloor t/P\rfloor\ge3\) and all \(d\) digits are nonzero, which
gives \(\operatorname{ord}(A)=h\).

In a distinct representation the exceptions contribute at most \(2S\),
where \[
\frac{S}{P}=\frac{1}{P}\sum_{c\in C} c
=1+\frac12\sum_{i=1}^d\frac{u_i-1}{P_{i-1}}.
\] The first digit contributes \(3/2\). For \(i\ge2\) we have
\(u_i/P_{i-1}=2^{2-i}\), so the remaining contribution is strictly less
than \(\frac12\sum_{i=2}^{\infty}2^{2-i}=1\). Thus \(2S<7P\). For
\(n\equiv8P-1\bmod 8P\), the exception sum plus the \textbf{number} of
residue-one tail terms is nonnegative and is congruent to \(8P-1\).
Hence at least \(8P-1-2S\ge P\) tail terms are required.

Finally, for each residue modulo \(m\) choose its representative
\(k\in\{1,\ldots,m\}\). Every sufficiently large integer \(n\) in that
class is a sum of \(k\) distinct residue-one tail members: fix \(k-1\)
distinct such members and use their complement in \(n\) for the last.
The same argument avoids any prescribed finite deletion. Therefore
\(\operatorname{ord}_R(A)\) exists and is at least
\(P=2^{2^{h-5}+h-4}\), proving Theorem 1. \(\square\)

\textbf{Finite checks.} As a separate verification, the accompanying
Python program computes both orders of the first four examples exactly.
Each modular representation must contain a tail member; conversely, any
fixed representation meeting this requirement lifts to every
sufficiently large target in its residue, also with distinct tail
members when required. Breadth-first search allows repetitions, while a
subset dynamic program allows each exception at most once and then fills
residues with additional tail ones.

\begin{longtable}[]{@{}
  >{\raggedleft\arraybackslash}p{(\columnwidth - 6\tabcolsep) * \real{0.2500}}
  >{\raggedleft\arraybackslash}p{(\columnwidth - 6\tabcolsep) * \real{0.2500}}
  >{\raggedleft\arraybackslash}p{(\columnwidth - 6\tabcolsep) * \real{0.2500}}
  >{\raggedleft\arraybackslash}p{(\columnwidth - 6\tabcolsep) * \real{0.2500}}@{}}
\toprule\noalign{}
\begin{minipage}[b]{\linewidth}\raggedleft
Ordinary order \(h\)
\end{minipage} & \begin{minipage}[b]{\linewidth}\raggedleft
Period \(m\)
\end{minipage} & \begin{minipage}[b]{\linewidth}\raggedleft
Sum of exceptions
\end{minipage} & \begin{minipage}[b]{\linewidth}\raggedleft
Exact restricted order
\end{minipage} \\
\midrule\noalign{}
\endhead
\bottomrule\noalign{}
\endlastfoot
5 & 32 & 20 & 15 \\
6 & 128 & 92 & 42 \\
7 & 1,024 & 792 & 245 \\
8 & 32,768 & 26,336 & 6,476 \\
\end{longtable}

\subsection{References}\label{references}

\begin{enumerate}
\def\labelenumi{\arabic{enumi}.}
\tightlist
\item
  P. Erdős and R. L. Graham.
  \href{https://www.renyi.hu/~p_erdos/1980-22.pdf}{On bases with an
  exact order}. \emph{Acta Arith.} \textbf{37} (1980), 201--207.
\item
  P. Erdős and R. L. Graham.
  \href{https://mathweb.ucsd.edu/~ronspubs/80_11_number_theory.pdf}{Old
  and New Problems and Results in Combinatorial Number Theory}.
  Monographies de L'Enseignement Mathématique \textbf{28} (1980).
\item
  N. Hegyvári, F. Hennecart and A. Plagne.
  \href{https://doi.org/10.1515/crll.2003.055}{A proof of two Erdős'
  conjectures on restricted addition and further results}. \emph{J.
  Reine Angew. Math.} \textbf{560} (2003), 199--220.
\item
  N. Hegyvári, F. Hennecart and A. Plagne.
  \href{https://www.cmls.polytechnique.fr/perso/plagne/Erdos-Burr.pdf}{Answer
  to a question by Burr and Erdős on restricted addition, and related
  results}. \emph{Combin. Probab. Comput.} \textbf{16} (2007), 747--756.
\item
  S.-Q. Chen, Q.-H. Yang and L. Zhao.
  \href{https://arxiv.org/html/2611.33191v1}{On a problem of Erdős and
  Graham related to restricted order}. arXiv:2611.33191v1 (2026).
  Preprint.
\item
  M. Kneser. \href{https://doi.org/10.1007/BF01174162}{Abschätzung der
  asymptotischen Dichte von Summenmengen}. \emph{Math. Z.} \textbf{58}
  (1953), 459--484.
\item
  N. Hegyvári.
  \href{https://real-d.mtak.hu/1038/7/dc_1260_16_doktori_mu.pdf}{Topics
  in Combinatorial Number Theory}. Doctoral dissertation, Hungarian
  Academy of Sciences (2017), Chapter 4, Theorem 4.8, p.~63.
  Quantitative form of the density theorem in {[}3{]}.
\end{enumerate}

\emph{Disclosure.} The finite disjoint-pair counting argument adapts
{[}5, Lemma 3{]}; the progression lemma is taken from {[}5{]}, not
claimed as a new result here. This draft was developed with AI
assistance. It has not been peer reviewed.

\end{document}
